\clearpage
\fancyhead[L]{\small MERCURSTAFELN.}
\SourceLeaf{564}
\begin{center}\small Göttingische gelehrte Anzeigen. Stück 107. Seite 1065..1069. 1813 Juli 5.\\[3pt]\rule{.72\linewidth}{.3pt}\end{center}

Herr \textsc{von Lindenau}, correspondirendes Mitglied der königl. Societät der Wissenschaften, welcher sich bereits durch seine Venus- und Marstafeln um die Planetentheorie verdient gemacht hat, ertheilte in einem Schreiben vom 20. Juni d. J. an den Professor \textsc{Gauss} eine vorläufige Nachricht von einer ähnlichen, nicht weniger wichtigen, dritten Arbeit über den \emph{Mercur}, welche er seit kurzem vollendet hat. Es haben zwar die verdienstlichen Bemühungen von \textsc{Lalande}, \textsc{Oriani} und \textsc{Triesnecker} die Theorie der Mercursbahn bereits zu einem hohen Grade von Vollkommenheit gebracht. Allein dessen ungeachtet war es wegen mancher Hülfsmittel, die, in practischer und theoretischer Hinsicht, theils nicht benutzt worden, theils neu hinzugekommen sind, nicht unwahrscheinlich, dass eine neue Bearbeitung dieses Gegenstandes, mit sorgfältiger Berücksichtigung alles dessen, was die heutige Astronomie zu diesem Behuf Dienliches darbietet, noch Aenderungen und Verbesserungen der vorhandenen Bestimmungen gewähren werde. Dieser Grund, verbunden mit der Hoffnung, durch eine genaue Vergleichung der durch die Theorie gegebenen Störungen der elliptischen Mercursbahn mit den aus den Beobachtungen folgenden eine neue und zuverlässige Bestimmung der Venusmasse herleiten zu können, war es hauptsächlich, was Hrn. \textsc{von Lindenau} zu seiner neuen Bearbeitung dieses Gegenstandes veranlasste. Die Details dieser Untersuchungen, wovon wir hier nur einige Endresultate beibringen, werden den Astronomen in einem, in wenig Wochen im Druck erscheinenden, Werke dargelegt werden, unter dem Titel: \emph{Investigatio nova orbitae a Mercurio circa Solem descriptae: accedunt tabulae Planetae ex elementis recens correctis et ex theoria gravitatis ill. de la Place constructae, auctore} \textsc{Bernhardo de Lindenau}.

Dem vorher angedeuteten Zwecke der Untersuchung gemäss, zerfällt diese in vier Abschnitte. In den beiden ersten werden die seit 1631 beobachteten Mercursdurchgänge, in den letztern neuere geocentrische Beobachtungen, discutirt und zur Bahnbestimmung benutzt. Ein doppeltes Verfahren ward auf die Mercursdurchgänge in Anwendung gebracht. Erstens konnte daraus Knoten, Knotenbewegung und Neigung bestimmt werden, und dann erlaubte zweitens das Eigenthümliche der Mercursbahn, vermöge dessen alle Durchgänge in einerlei Knoten nahe in denselben Punkten der Bahn Statt finden, die durch die Durchgänge gegebenen heliocentrischen Längen in der Ekliptik auf mittlere in der Bahn zu reduciren, ohne dass die Differenzen zweier durch andere Elemente, als durch die mittlere jährliche Bewegung und die jährliche Aenderung des Aphelium afficirt und wesentlich irrig gemacht werden könnten. Doch muss man zu Anwendung dieses Verfahrens allerdings eine schon genäherte Kenntniss der Mercursbahn besitzen, wie dies denn jetzt wirklich der Fall war. Zu der erstern Bestimmung, die vorzüglich eine genaue Kenntniss der geocentrischen Breite erforderte, konnten nur zwölf Durchgänge benutzt werden, und aus den siebenzehn, die überhaupt von dem Verfasser discutirt und berechnet worden sind, wurden funfzehn Combinationen gebildet, die ihm zu numerischer Entwickelung und Anwendung des zweiten Verfahrens am vortheilhaftesten geeignet schienen. Die Vergleichung der hieraus erhaltenen jährlichen Aenderungen des Knotens und der Sonnenferne mit den durch die Theorie gegebenen, deren Werthe hauptsächlich von der Venusmasse abhängen, gab die beiden ersten Gleichungen zu deren Bestimmung an die Hand.

Die aus den nur erwähnten siebenzehn Durchgängen hergeleiteten heliocentrischen Mercurslängen, nebst hundert von \textsc{Maskelyne} und \textsc{Piazzi} beobachteten geocentrischen Längen, dienten zur Bestimmung der eigentlich elliptischen Elemente der Mercursbahn. Da für Neigung und Knoten schon vorher genäherte Werthe erhalten worden waren, so konnten diese, auf die Reduction nur geringen \SourceLeaf{565}Einfluss habenden, Elemente zuerst ohne Bedenken unberücksichtigt bleiben. Die Differenzen der beobachteten und berechneten Längen wurden durch eine Function der Correctionen der Epoche, der mittlern Bewegung, der Excentricität, der Sonnenferne und der Venusmasse ausgedrückt, und so 117 Bedingungsgleichungen formirt, aus denen die Correctionen der zum Grunde gelegten Elemente nach der Methode der kleinsten Quadrate hergeleitet wurden. Daraus ergab sich dann auch die dritte Gleichung für den Werth der Venusmasse. Die Formation der Bedingungsgleichungen aus den geocentrischen Längen erforderte Vorsicht und Schärfe, indem hier, wo der heliocentrische Fehler in den grössten Elongationen (in denen Mercur hauptsächlich und fast einzig beobachtet ist), durch Reduction auf den geocentrischen Ort sehr vermindert wird, kleine Fehler in den beobachteten geocentrischen Längen die gesuchten Correctionen der Elemente stark ändern. Diesen Einfluss der Beobachtungsfehler glaubt der Verf. durch die Anzahl der Beobachtungen eliminirt, und den der Sonnenörter auf die berechneten Längen durch jedesmalige Verbesserung der Sonnentafeln aus gleichzeitigen Beobachtungen, vermieden zu haben.

Mit diesen verbesserten Elementen wurden die beobachteten geocentrischen Breiten auf heliocentrische reducirt, und aus deren Vergleichung mit den berechneten hundert Bedingungsgleichungen formirt, welche Neigung und Knoten gaben. Der Wunsch, die Neigung für eine frühere Epoche zu bestimmen, und sonach aus deren beobachteter Säcularänderung noch eine vierte Gleichung für die Venusmasse zu erhalten, wurde durch Mangel tauglicher Beobachtungen vereitelt.

Die Endresultate dieser Untersuchungen waren folgende:

Wird die Venusmasse, wie sie \textsc{Laplace} (Mécan. Céleste T.\,3. p.\,61) annimmt, $= 1$ gesetzt, so ist die verbesserte, wie sie aus den erwähnten drei Gleichungen folgt, $= 1{,}0974 = \frac{1}{849132}$ der Sonnenmasse. Mit dieser Masse sind alle periodischen und Säcularstörungen berechnet worden.

\begin{center}\fontsize{9}{11}\selectfont\setlength{\tabcolsep}{2pt}\renewcommand{\arraystretch}{1.18}\begin{tabular}{@{}l c r@{\;}r@{\;}r@{}}
Epoche 1750 Meridian von Seeberg & $.\;.\;.\;.\;.\;.$ & $253^\circ$ & $5'$ & $17''1$\\
Mittlere jährliche Bewegung & $.\;.\;.\;.\;.\;.\;.\;.$ & $53$ & $43$ & $3.613$\\
Sonnenferne 1750 & $.\;.\;.\;.\;.\;.\;.\;.\;.\;.\;.\;.\;.$ & $253$ & $33$ & $24.3$\\
Excentricität 1800 & $.\;.\;.\;.\;.\;.\;.\;.\;.\;.\;.\;.$ & \multicolumn{3}{l}{$0.2056163$}\\
Halbe grosse Axe & $.\;.\;.\;.\;.\;.\;.\;.\;.\;.\;.\;.$ & \multicolumn{3}{l}{$0.3870988$}\\
Knoten 1750 & $.\;.\;.\;.\;.\;.\;.\;.\;.\;.\;.\;.\;.\;.\;.$ & $75^\circ$ & $22'$ & $0''96$\\
Neigung der Bahn 1800 & $.\;.\;.\;.\;.\;.\;.\;.\;.\;.$ & $7$ & $0$ & $5.9$\\
Säcularänderung der Sonnenferne & $.\;.\;.\;.\;.\;.\;.$ & $+1$ & $33$ & $22.9$\\
\quad —\quad —\quad des Knoten & $.\;.\;.\;.\;.\;.\;.\;.$ & $+1$ & $10$ & $15.1$\\
\quad —\quad —\quad der Excentricität & $.\;.\;.\;.\;.\;.$ & \multicolumn{3}{l}{$+0''791$}\\
\quad —\quad —\quad der Neigung & $.\;.\;.\;.\;.\;.\;.\;.$ & \multicolumn{3}{l}{$+18.380$}\\
\end{tabular}\end{center}

{\rightskip=0pt plus 3em Sämmtliche hundert in die Bedingungsgleichungen aufgenommene beobachtete geocentrische Oerter werden durch diese Elemente äusserst befriedigend dargestellt, und eben dies ist bei funfzig neuern, von \textsc{Bouvard} in Paris und vom Verf. selbst auf der Seeberger Sternwarte angestellten Mercursbeobachtungen, die nicht mit zur Begründung der Elemente dienten, der Fall, so dass man zu der Hoffnung vollkommen berechtigt ist, dass diese Bestimmungen die Mercursbewegungen auch in den nächsten Jahrzehnten mit dem Himmel übereinstimmend darstellen werden.\par}
